
To derive the dynamical model of the vehicle, we start by modeling the
force and torque vectors caused by a single propeller in the body
frame denoted~$\mathcal{B}$, whose origin is coincident with the
vehicle's center of mass (CoM).

Consider a propeller $i$ whose motor is rigidly linked to the
body frame, as depicted in Figure~\ref{fig:prop}.  This propeller will
generate a force $\bar{F}_i$ and a torque $\bar{M}_i$ with respect to
the CoM. While the former results directly from the
propeller thrust,
\begin{align}
  \label{eq:Fi}
  \bar{F}_i &= f_i \hat{u}_i \\
  f_i &= K_1 u_i
\end{align}
where $u_i$ is the actuation signal and $f_i$ is the scalar thrust,
the later has two components: the torque caused by the displaced force
and the propeller reaction torque:
\begin{align}
  \label{eq:Mi}
  \bar{M}_i &= \bar{r}_i\times\bar{F}_i - \tau_i \hat{u}_i \\
  \tau_i &= w_i K_2 u_i
\end{align}
where $\times$ denotes vector cross product and $w_i$ is $-1$ or $1$ depending on whether the propeller rotates
clockwise or anti-clockwise for a positive forward thrust $f_i>0$.
The constants $K_1$ and $K_2$ model a simple linear relationship
between the actuation and the scalar thrust force and reaction torque
of the propeller. The actuation is often modeled as proportional to
the square of the propeller speed of rotation~\cite{mahony12}.

The relative position of the propeller, $\bar{r}_i$, orthogonal to the
$Z$ axis, and the unit vector $\hat{u}_i$, aligned with the propeller
axis, are uniquely defined by the angles $\theta_i$ and $\phi_i$, and can be
easily obtained from geometrical reasoning:
\begin{equation}
  \bar{r}_i = \left(
    \begin{array}{c}
      d\cos(\theta_i) \\
      d\sin(\theta_i) \\
      0
      \end{array}
    \right)
    \qquad
    \hat{u}_i = \left(
      \begin{array}{c}
        \sin(\theta_i)\sin(\phi_i) \\
        -\cos(\theta_i)\sin(\phi_i) \\
        \cos(\phi_i)
      \end{array}
      \right)
\end{equation}
where $d=\|\bar{r}_i\|$ is the distance from the propeller to the CoM.
Stacking together the resulting force $\bar{F}_i$ and torque
$\bar{M}_i$, one can obtain the linear relation
\begin{equation}
  \label{eq:fm}
  \left(
    \begin{array}{c}
      \bar{F}_i \\
      \bar{M}_i
    \end{array}
  \right) = \bar{a}_i\,u_i
\end{equation}
where
\begin{equation}
  \begin{split}
    \bar{a}_i  &= \left(
      \begin{array}{c}
        K_1 \hat{u}_i \\
        K_1 \bar{r}_i \times \hat{u}_i - w_i K_2 \hat{u}_i
      \end{array}
    \right) \\
    &= \left(
      \begin{array}{c}
        K_1 \sin(\theta_i)\sin(\phi_i) \\
        -K_1 \cos(\theta_i)\sin(\phi_i) \\
        K_1 \cos(\phi_i) \\
        (K_1 d \cos(\phi_i) - w_i K_2 \sin(\phi_i))\sin(\theta_i) \\
        -(K_1 d \cos(\phi_i) - w_i K_2 \sin(\phi_i))\cos(\theta_i) \\
        -K_1 d \sin(\phi_i) - w_i K_2 \cos(\phi_i)
      \end{array}
    \right)
    \end{split}
\end{equation}

\begin{figure}
  \centering
  \includegraphics[width=0.6\linewidth]{prop.pdf}
  \caption{Notation used for modeling a single propeller with respect
    to the body frame~$\mathcal{B}$, centered on the vehicle's CoM.}
  \label{fig:prop}
\end{figure}

For a vehicle with $N$ propellers, the resulting force and torque will
be given by the sum of the contributions of each propeller, given with
respect to the CoM, by~\eqref{eq:fm}. This sum can be put in matrix
form as
\begin{equation}
  \label{eq:au}
  \left(
    \begin{array}{c}
      \bar{F} \\
      \bar{M}
    \end{array}
  \right) = \mathbf{A}\,\bar{u}
\end{equation}
where $\mathbf{A}=[\bar{a}_1 \cdots \bar{a}_N]$, hereby called
\emph{actuation matrix}, and $\bar{u}=[u_1\cdots u_N]^T$, the
actuation input vector. The crucial observation is that, if the
actuation matrix $\mathbf{A}$ has at least rank 6, the six degrees of
freedom of the force and torque combined can be decoupled by
solving~\eqref{eq:au} with respect to the actuation~$\bar{u}$. A
necessary (but not sufficient) condition for this to be true is to have at least 6
propellers, as in the hexarotor design considered in this paper. Note
that this matrix only depends on the vehicle's design parameters,
defined by the angles $\{\theta_i\}$ and $\{\phi_i\}$, distance $d$,
and the trust coefficients $K_1$, $K_2$, and $\{w_i\}$. However, we
will defer the discussion about the properties of this matrix to the
next section.

The dynamical model of the vehicle, in the absence of gravity can be
derived from the Newton and Euler equations of motion, similarly
to~\cite{mahony12}. Let us denote the position and velocity of the
body frame $\mathcal{B}$ with respect to the inertial frame
$\mathcal{I}$ as $\bar{x}$ and $\bar{v}$, the rotation matrix of frame
$\mathcal{B}$ with respect to $\mathcal{I}$ as $\mathbf{R}$, and the
angular velocity of the vehicle in the body frame $\mathcal{B}$ as
$\bar{\omega}$. Then,
\begin{equation}
  \label{eq:dynsys}
  \left\{
    \begin{aligned}
      \dot{\bar{x}} &= \bar{v} \\
      m \dot{\bar{v}} &= \mathbf{R} \bar{F} \\
      \dot{\mathbf{R}} &= \mathbf{R}\,\mathbf{S}(\bar{\omega}) \\
      \mathbf{J} \dot{\bar{\omega}} &= \bar{M} - \bar{\omega}\times\mathbf{J} \bar{\omega}
    \end{aligned}
  \right.
\end{equation}
where the constants $m$ and $\mathbf{J}$
are the vehicle's mass and moment of inertia, while
$\mathbf{S}(\bar{\omega})$ is the skew-symmetric matrix defined by
\begin{equation}
  \label{eq:S}
  \mathbf{S}(\bar{\omega}) =
  \left[
    \begin{array}{ccc}
      0 & -\omega_z & \omega_y \\
      \omega_z & 0 & -\omega_x \\
      -\omega_y & \omega_x & 0 \\
    \end{array}
  \right]
\end{equation}
for $\bar{\omega}=[\omega_x\:\omega_y\:\omega_z]^T$.

%%% Local Variables:
%%% mode: latex
%%% TeX-master: "main"
%%% End:
